% ---------------------------------------------------------------
\section{An Assurance Model for Training-Data Pipelines}\label{sec:model}
% ---------------------------------------------------------------

This section turns the informal question ``what does a provenance record prove?''
into a decision problem. We fix a pipeline shape, an adversary, a set of mechanisms and
an auditor's acceptance rule, and ask for which combinations an \emph{accepted}
execution is guaranteed to satisfy a goal.

\subsection{Pipeline and ground truth}

A training run passes data through five stages (Fig.~\ref{fig:pipeline}):
\emph{authorised sources} deliver items; an optional \emph{generator} produces synthetic
items; a \emph{transform} (filtering, deduplication, mixing) produces the training input;
a \emph{capture} stage commits a record of that input; and an \emph{executor} trains and
releases a model. Items are atomic and content-identified: two different contents are two
different items.

Let $\mathcal{I}$ be a finite universe of items. An execution fixes eleven subsets of
$\mathcal{I}$, which describe what actually happened (ground truth) and what the auditor
is told (reports):
\begin{center}
\begin{tabular}{@{}ll@{}}
\toprule
$D$ & items delivered by authorised sources\\
$S$ & items produced by the declared generator\\
$F$ & transform output = items delivered to the executor\\
$C$ & items committed in the record\\
$U$ & items that determine the released model\\
$B$ & poisoned items\\
$M\subseteq B$ & poisoned items a detector would miss\\
$X$ & items an attested detector examined\\
$\hat S$ & synthetic set as reported to the auditor\\
$\hat D$ & deliveries as reported by source witnesses\\
$\hat F$ & training inputs as reported by an input witness\\
\bottomrule
\end{tabular}
\end{center}
Two global Booleans complete an execution: $\mathit{gen}$ (a generative stage is declared)
and $\mathit{gs}$ (the training seed was chosen by the prover after seeing the data).

\subsection{Adversary}

The adversary is a set $\alpha$ of capabilities, one per point of trust:
\textsf{AdvSource} (an authorised source delivers poisoned data),
\textsf{AdvGen} (the generator emits poisoned data),
\textsf{AdvTrans} (the transform operator is corrupt),
\textsf{AdvCap} (the capture operator is corrupt),
\textsf{AdvExec} (the training executor is corrupt),
\textsf{AdvWit} (witnesses collude) and
\textsf{AdvTEE} (attestation keys are extracted or the enclave is broken).
We study six profiles (Table~\ref{tab:profiles}). The \emph{full} adversary controls every
pipeline stage but not the trust base; \emph{full+wit} and \emph{full+tee} additionally break
one trust anchor each; the three single-stage profiles model deployments in which only one
operator is untrusted. Every profile includes \textsf{AdvSource}: an authorised source can
always be wrong, because authorisation is a statement about \emph{who} supplied data, not
about \emph{what} the data contains.

The six profiles are \emph{illustrative}, chosen to match common deployments. They are not
exhaustive. Section~\ref{sec:results} therefore also decides the model for all $2^7=128$
subsets of capabilities. A subset that combines capabilities models \emph{correlated
compromise}. For example, \textsf{AdvWit} together with \textsf{AdvTEE} is a single
administrator who controls both the witness service and the attestation keys. The model
does not separate failures \emph{within} a mechanism class: \textsf{AdvWit} corrupts source
and input witnesses together, and \textsf{AdvTEE} breaks every enclave. Finer correlation
structures, such as one compromised enclave among several, are outside the model.

\paragraph{Trusted parties}
The auditor, its verification code, the public randomness beacon used by
\textsf{BindFix}, and the cryptographic primitives are trusted. The record store is assumed
to provide integrity and non-equivocation ($P_1$, $P_2$; Proposition~\ref{prop:ledger}).
Everything else may be corrupted according to $\alpha$.

\begin{table}[t]
\centering\small
\caption{Adversary profiles.}\label{tab:profiles}
\begin{tabular}{@{}lccccccc@{}}
\toprule
Profile & Src & Gen & Trans & Cap & Exec & Wit & TEE\\
\midrule
full        & \checkmark & \checkmark & \checkmark & \checkmark & \checkmark & & \\
full+wit    & \checkmark & \checkmark & \checkmark & \checkmark & \checkmark & \checkmark & \\
full+tee    & \checkmark & \checkmark & \checkmark & \checkmark & \checkmark & & \checkmark\\
cap-only    & \checkmark & & & \checkmark & & & \\
exec-only   & \checkmark & & & & \checkmark & & \\
trans-only  & \checkmark & \checkmark & \checkmark & & & & \\
\bottomrule
\end{tabular}
\end{table}

\subsection{Mechanisms}

The mechanism set $\mu$ is a subset of seven primitives, each an abstraction of a family
of deployed techniques (Section~\ref{sec:sok}):
\begin{itemize}
\item \textsf{BindFix}, \textsf{BindGrind}: \emph{computation binding}, i.e.\ the auditor
  can confirm that the released model was produced from the committed record (re-execution,
  proof of learning, zero-knowledge proof of training). \textsf{BindFix} fixes the
  randomness independently of the prover (e.g.\ from a public beacon released after
  commitment); \textsf{BindGrind} accepts a prover-chosen seed.
\item \textsf{WitSrc}: \emph{source witnessing}, i.e.\ authorised sources sign
  content-addressed delivery receipts, and the auditor rejects committed items covered
  neither by a receipt nor by the synthetic report.
\item \textsf{WitIn}: \emph{input witnessing}, i.e.\ an independent observer under complete
  mediation reports the items delivered to the executor, and the auditor rejects a record
  that omits one.
\item \textsf{AttCap}, \textsf{AttTrans}: \emph{attested capture} and \emph{attested
  transform}, i.e.\ the stage runs in a trusted execution environment whose signed
  attestation binds its input to its output.
\item \textsf{Detector}: an attested content detector is run over the committed record.
\end{itemize}
Write $\beta\equiv \textsf{BindFix}\in\mu\vee\textsf{BindGrind}\in\mu$ and
$\tau\equiv\textsf{AdvTEE}\notin\alpha$.

\subsection{Semantics}

An execution is \emph{admissible} for $(\alpha,\mu,\mathit{gen})$ if it satisfies the
constraints $\Phi$ below. Each constraint states what a stage does when it is honest, or
when a mechanism forces it to behave honestly.
{\small
\begin{align*}
&(\Phi_1)\ D\cap S=\emptyset;\quad \neg\mathit{gen}\Rightarrow S=\hat S=\emptyset;\quad M\subseteq B\\
&(\Phi_2)\ \textsf{AdvSource}\notin\alpha\Rightarrow B\cap D=\emptyset\\
&\phantom{(\Phi_2)}\ \textsf{AdvGen}\notin\alpha\Rightarrow B\cap S=\emptyset\\
&(\Phi_3)\ \textsf{AdvTrans}\notin\alpha \vee (\textsf{AttTrans}\in\mu\wedge\tau)\Rightarrow F\subseteq D\cup S\\
&(\Phi_4)\ (\textsf{AttTrans}\in\mu\wedge\tau)\vee(\textsf{AdvTrans},\textsf{AdvCap}\notin\alpha)\\
&\phantom{(\Phi_4)}\ \Rightarrow \hat S=S\\
&(\Phi_5)\ \textsf{AdvCap}\notin\alpha\vee(\textsf{AttCap}\in\mu\wedge\tau)\Rightarrow C=F\\
&(\Phi_6)\ \textsf{AdvExec}\notin\alpha\Rightarrow U=F\\
&\phantom{(\Phi_6)}\ \mathit{gs}\Rightarrow\textsf{AdvExec}\in\alpha\wedge\textsf{BindFix}\notin\mu\\
&(\Phi_7)\ \textsf{AdvWit}\notin\alpha\Rightarrow\hat D=D\wedge \hat F=F\\
&(\Phi_8)\ \textsf{Detector}\in\mu\wedge\tau\Rightarrow X=C;\ \text{otherwise } X=\emptyset
\end{align*}}
The second clause of $\Phi_1$ encodes \emph{declared stage types}: if no generator is
declared, nothing can be reported as synthetic. Section~\ref{sec:results} shows that this is
a necessary design requirement, not a modelling convenience.

The auditor \emph{accepts} an execution if
{\small
\begin{align*}
&(\mathrm{A}_1)\ \beta\Rightarrow U=C\\
&(\mathrm{A}_2)\ \textsf{WitSrc}\in\mu\Rightarrow C\subseteq \hat D\cup\hat S\\
&(\mathrm{A}_3)\ \textsf{WitIn}\in\mu\Rightarrow \hat F\subseteq C\\
&(\mathrm{A}_4)\ \textsf{Detector}\in\mu\wedge\tau\Rightarrow C\cap B\subseteq M
\end{align*}}
$(\mathrm{A}_4)$ says that a detected poisoned item causes rejection; $M$ is the set the
detector fails to flag, so an imperfect detector is any $M\neq\emptyset$.

\subsection{Assurance goals}

\begin{center}\small
\begin{tabular}{@{}llp{3.3cm}@{}}
\toprule
Goal & Condition & Reading\\
\midrule
$G_{\mathrm{bind}}$ & $U=C$ & the record explains the model's inputs\\
$G_{\mathrm{auth}}$ & $U\subseteq D\cup S$ & no unauthorised lineage\\
$G_{\mathrm{noinf}}$ & $U=C\wedge\neg\mathit{gs}$ & nothing outside the record shaped the model\\
$G_{\mathrm{content}}$ & $U\cap B\subseteq M\cap X$ & residual poison was examined and missed\\
$G_{\mathrm{zero}}$ & $U\cap B=\emptyset$ & no poison was used\\
$G_{\mathrm{compl}}$ & $F\subseteq C$ & the record omits no training input\\
\bottomrule
\end{tabular}
\end{center}

\begin{definition}[Assurance]
Mechanism set $\mu$ \emph{achieves} goal $G$ against $\alpha$ for pipeline type
$\mathit{gen}$, written $\mu\models_{\alpha,\mathit{gen}}G$, if every admissible and accepted
execution, over every finite universe $\mathcal{I}$, satisfies $G$.
\end{definition}

The definition quantifies over all universes, so it is not a bounded statement. The next
section shows that it can nevertheless be decided exactly.

\subsection{Scope of the abstraction}\label{sec:scope}

The model treats items as atomic, content-identified units. Each mechanism either is present
and enforced, or absent. A single training run is modelled. These choices make the results
exact, and they also bound what is proved. Table~\ref{tab:scope} states how common LLM
training practices relate to the abstraction.

\begin{table}[t]
\centering\scriptsize
\caption{How LLM training practices relate to the model's abstraction.}\label{tab:scope}
\begin{tabular}{@{}p{2.2cm}p{5.8cm}@{}}
\toprule
Practice & Status in the model\\
\midrule
Tokenisation, document segmentation, sequence packing & Covered if the \emph{item} is the
unit that is committed and receipted (e.g.\ a document or a packed sequence). The results apply
at that granularity. Claims at finer granularity need receipts and commitments at that
granularity.\\
Partial inclusion (truncation, sampling) & A transform. Covered when the truncated or sampled
item is itself the committed item; the original is then a different item.\\
Data mixing, deduplication, filtering & The transform stage ($F$).\\
Synthetic data generation & The generator stage ($S$), with \textsf{AdvGen} and declared stage
types.\\
Distributed or data-parallel training & One executor stage. Binding must cover all workers;
partial binding of some workers is not modelled.\\
Checkpoint selection & A prover choice among accepted outputs, i.e.\ grinding
(Theorem~\ref{thm:grind}) unless the selection rule is fixed in advance.\\
Successive fine-tuning, distillation, model merging & Outside the single-run model. A chain
of runs is covered only if each run is modelled separately and the previous model is
treated as an input item (see the remark below).\\
Intermediate representations, gradients & Not modelled.\\
\bottomrule
\end{tabular}
\end{table}

\paragraph{Remark (multi-stage composition, argued)}
Consider a chain of runs in which run $j$ trains on record $C_j$ and the released model
$M_{j-1}$ of run $j-1$. Treat $M_{j-1}$ as one item of $C_j$, authorised only if run $j-1$
achieved the same goal. Then a goal holds for the final model if it holds for every run.
For $G_{\mathrm{bind}}$, $G_{\mathrm{auth}}$ and $G_{\mathrm{noinf}}$ this is immediate by
induction. For $G_{\mathrm{content}}$ it is not: a detector applied to $C_j$ does not
examine the data that shaped $M_{j-1}$. So content claims do not compose without a detector
at every stage. We state this composition as an argued remark, not as a theorem. Model
merging and distillation from several parents fall outside it.
